\documentclass[10pt,a4paper]{article}

\usepackage[margin=1.9cm]{geometry}
\usepackage{booktabs}
\usepackage{microtype}
\usepackage{times}
\usepackage{enumitem}

\setlist{nosep,leftmargin=*}
\setlength{\parskip}{3pt}
\setlength{\parindent}{0pt}

\title{\vspace{-1.4cm}\textbf{Project Plan --- P02: Plant Disease Recognition from Leaf Images}\vspace{-0.2cm}}
\author{\small Egor Savchenko \quad Zamir Safin \quad Ilia Ponomarev \quad Vadim Poponnikov\\[2pt]
\small Innopolis University --- Computer Vision 2026\vspace{-0.3cm}}
\date{\small 6 October 2026}

\begin{document}
\maketitle

\section*{1. Project question}

\textbf{Main question.} Does full fine-tuning of an ImageNet-pretrained ResNet-50 retain its
advantage over frozen-backbone linear probing when leaf images are degraded by realistic robotic
field-capture artifacts rather than captured under laboratory conditions?

The controlled variable is which parameters receive gradients. Architecture, data, optimiser,
schedule and seed are held fixed, so any difference between the two conditions is attributable to the
transfer strategy alone.

\textbf{Sub-questions.}
\begin{itemize}
  \item What is the clean-data gap between the two conditions, and at what training cost?
  \item How sensitive is the fine-tuned condition to the learning rate?
  \item How does each condition degrade under three independent capture artifacts at three severities?
  \item Which visually similar disease pairs are confused, and which classes fail first under
        degradation? (Required explicitly by the P02 brief.)
\end{itemize}

\section*{2. Dataset}

PlantVillage, Kaggle mirror \texttt{mohitsingh1804/plantvillage}: 38 classes across 14 crops,
laboratory-captured images of single detached leaves on a uniform background, $256\times256$ px.
Chosen because it is public, class-balanced enough for macro-averaged metrics, and small enough that
the whole study fits in the course scope without large-scale training.

Preprocessing: directory names are sanitised, images are resized to $224\times224$ and normalised
with ImageNet statistics. Training augmentation is horizontal and vertical flips plus rotation up to
$15^{\circ}$; validation and test use no augmentation.

\section*{3. Train / validation / test split}

\begin{tabular}{lrl}
\toprule
\textbf{Split} & \textbf{Images} & \textbf{Origin} \\
\midrule
Train & 41\,275 & Kaggle \texttt{train}, minus the test sample \\
Validation & 10\,861 & Kaggle \texttt{val}, used as shipped \\
Test & 2\,169 & 5\,\% of each training class, seed 42 \\
\bottomrule
\end{tabular}

\textbf{Leakage control.} The test split is built by \emph{moving} files out of \texttt{train}, never
by copying, so no image can appear in two splits. The split script is deterministic under seed 42 and
is run exactly once on a fresh download. Class lists are asserted identical across the three splits
before training starts. Model selection uses validation macro F1 only; the test split is read once
per run, after training has finished.

\section*{4. Evaluation measures}

\begin{itemize}
  \item \textbf{Accuracy} --- headline measure, directly comparable with published PlantVillage results.
  \item \textbf{Macro F1} --- primary measure for model selection and for the comparison, because it
        weights all 38 classes equally and is not dominated by the largest ones.
  \item \textbf{Per-class precision, recall and F1} --- to locate which diseases fail.
  \item \textbf{Confusion pairs} --- off-diagonal mass of the confusion matrix, to answer the
        visually-similar-disease question in the brief.
  \item \textbf{Mean confidence on correct and on incorrect predictions} --- calibration under
        distribution shift.
  \item \textbf{Cost} --- trainable parameters, epochs to early stop, training time, inference
        latency at batch 1 and 64, and peak GPU memory.
\end{itemize}

\section*{5. Baseline and second condition}

\textbf{Baseline:} ResNet-50 with ImageNet weights, backbone frozen, only the 38-way classifier
trained (77\,862 of 23\,585\,894 parameters) --- a linear probe on fixed features.

\textbf{Second condition:} the same network with all 23\,585\,894 parameters trainable.

\textbf{Ablation:} learning rate $10^{-4}$, $10^{-3}$, $10^{-2}$ on the fine-tuned condition.

\textbf{Degradation study:} both trained models are evaluated, without retraining, on the clean test
set and on nine degraded versions of it, modelling an autonomous field phenotyping robot ---
directional motion blur, low light with sensor noise and auto-gain, and JPEG compression, each at
three severities.

\section*{6. Division of work}

\begin{tabular}{@{}p{3.1cm}p{13.6cm}@{}}
\toprule
\textbf{Member} & \textbf{Responsibility} \\
\midrule
Egor Savchenko &
Overall architecture of the study and of the codebase. Formulation of the research question and of
the experimental design. Implementation of the full pipeline: training in both conditions, the
learning-rate ablation, the capture-degradation model, the inference and latency benchmark, the
failure and confusion analysis, and the results assembly. Testing and verification: recomputation of
every reported metric from the raw per-image predictions, cross-checking of independent runs, and the
reproducibility guarantees. \\[4pt]
Zamir Safin &
Dataset analysis and data handling: selection of the dataset and verification of its source and
licence; inventory of the 38 classes across 14 crops and their per-class counts; construction of the
held-out test split and verification that no image appears in two splits; sanitisation of directory
and file names; verification that the class lists of the three splits are identical; analysis of
class balance and identification of under-represented classes; the preprocessing and augmentation
policy; and the dataset provenance and download instructions in the README. \\[4pt]
Ilia Ponomarev &
Analysis of existing solutions and architectures: survey of transfer-learning approaches used for
plant disease recognition; justification of ImageNet-pretrained ResNet-50 as the backbone; framing of
frozen-backbone probing versus full fine-tuning as the controlled comparison; review of corruption
and robustness benchmarking practice, which motivates the capture-degradation design; and
positioning of our numbers against published PlantVillage baselines. \\[4pt]
Vadim Poponnikov &
Reporting, demo and presentation: the two-page technical summary, the presentation slides, and the
live demo including its scenario, the failure case shown, and rehearsal against the time limit. \\
\bottomrule
\end{tabular}

\section*{7. Schedule}

\begin{tabular}{@{}llp{10.4cm}@{}}
\toprule
\textbf{Date} & \textbf{Milestone} & \textbf{Deliverable} \\
\midrule
6 Oct & Project plan & This document \\
13 Oct & First working model & Dataset pipeline and baseline results \\
20 Oct & Second approach & Fine-tuned condition and the comparison \\
27 Oct & Progress check & Preliminary comparison and remaining experiments \\
3 Nov & Results checkpoint & Final metrics, runtime, tables and figures, failure cases \\
10 Nov & Final submission & Two-page summary, repository, slides, demo, contributions \\
\bottomrule
\end{tabular}

\end{document}
